\Section{Using Function Values in Move Specifications}
\label{sec:move-hof}

This section recalls the aspects of Move on Aptos \cite{MoveOnAptosBook} that matter for the verification of higher-order programs, introduces an automated market maker (AMM) as a running example, and illustrates the common specification patterns our extension supports.
All examples can be also found at \cite{PaperExamples}.

\SubSection{Move on Aptos} \label{sec:move}

Move is a bytecode-verifiable programming language designed from the outset to coexist with formal verification.
For the purposes of this paper, the relevant properties are:

\begin{Itemize}
    \item \emph{Strong static typing with generics.}  Every expression has a statically known explicit type; generic functions and structs are monomorphized for verification~\cite{DillGPQXZ22}.


    \item \emph{Type-indexed global storage.}  Global state is a map keyed by (resource type, address) pairs: a struct can be moved to, borrowed, and moved from the location !T[address]!, with linear ownership enforced statically.  Specifications quantify over this memory directly, and !modifies! clauses name the footprint a function may touch, providing a frame condition.

    \item \emph{Reference semantics.}  References are borrowed from local variables or from global storage; this yields an alias-free memory model that makes verification tractable~\cite{DillGPQXZ22}.

    \item \emph{Strongly-typed bytecode.}  Move source is compiled to bytecode whose type- and memory-safety are checked before execution.  \MVP verifies this bytecode against specifications in the source.

    \item \emph{Integrated specification language.}  Move lets developers write pre-/post-conditions, abort conditions, data invariants, and global memory invariants in the same source file as the code they describe, or optionally,
    into separate files.
\end{Itemize}


\Paragraph{The Move Prover}
\MVP is an automated formal verifier built into the Move toolchain;
the version discussed here and described in~\cite{DillGPQXZ22} is a complete rewrite of the earlier version in~\cite{MOVE_PROVER}.
Developers annotate Move source with specifications --- preconditions, postconditions, abort conditions, data invariants, and loop invariants --- using a first-order specification language that shares Move's type system.
\MVP translates the annotated program to Boogie~\cite{boogie,thisisboogie2}, an intermediate verification language, which Z3~\cite{z3} then discharges as SMT queries.
The tool targets routine use during development: verification is designed to complete within seconds per function through careful SMT encoding and monomorphized, modular function summaries; it is used in CI (continuous integration testing) in the workflows at Aptos.
Abort conditions deserve special note: multiple !aborts_if! clauses form a \emph{biconditional} --- the function aborts if and only if at least one clause holds --- rather than the one-directional ``if it aborts, then ...'' familiar from partial-correctness Hoare logic.

\Paragraph{Function values in Move}
A new version of Move extends the type system with \emph{function types} $T_1 \times \dots \times T_n \to T$ (written !|T1,...,Tn|T! in the concrete syntax) and allows values of such types to be constructed from concrete function symbols or from lambda expressions, possibly with captured arguments.
In Move's type ability system, function values may be marked !copy!, !drop!, and !store!; a function value with !store! can be persisted to global storage.
References may not be captured; only plain values may appear in a closure's payload, ensuring that captured data is independent of any borrow scope.

\SubSection{Automated Market Maker Example}
\label{sec:amm}

The running example throughout this paper is a simplified AMM module (Automated Market Maker, also called Uniswap in the Ethereum context~\cite{UniswapV2}) in which the \emph{pricing curve} is a function value stored in the AMM pool.
The pricing function has to satisfy a number of invariants: (i)~the function never aborts, (ii)~the output never exceeds the output reserve, (iii)~the output is monotone in the input amount, and (iv)~the product of the two reserves does not decrease after a swap (constant-product preservation).
The spec expresses these invariants using \emph{behavioral predicates} and \emph{state labels}, two constructs introduced by this paper:

% @formatter:off
\begin{MoveBox}
  struct Pool has key {
    reserve_x: u64, reserve_y: u64,
    // (reserve_in, reserve_out, amt) = out
    pricing: |u64, u64, u64|u64 has copy, store
  }
  spec Pool {
    // Function allowed to read any global state.
    reads_of<self.pricing> *;

    // Pricing must never abort.
    invariant forall S in *, ri: u64, ro: u64, a: u64:
      S |~ !aborts_of<self.pricing>(ri, ro, a);

    // Output never exceeds the output reserve.
    invariant forall S in *, ri: u64, ro: u64, a: u64:
     S.. |~ result_of<self.pricing>(ri, ro, a) <= ro;

    // Monotonicity in the input amount.
    invariant forall S in *, ri: u64, ro: u64,
                        a1: u64, a2: u64:
      S.. |~ a1 <= a2 ==>
        result_of<self.pricing>(ri, ro, a1)
             <= result_of<self.pricing>(ri, ro, a2);

    // Constant-product preservation.
    invariant forall S in *, ri: u64, ro: u64, a: u64:
      S.. |~
        (ri + a)
        * (ro - result_of<self.pricing>(ri, ro, a))
          >= ri * ro;
  }
\end{MoveBox}
% @formatter:on

Each invariant is universally quantified over a state label !S!, ranging over all states reachable during execution in which the pool exists at some address.
A \emph{state label} such as !S! names a program point; the state at !S! comprises both the global memory state and the local state (via references) at that point.
For the Pool invariants the local component is not relevant, since !pricing! takes only plain value arguments and no references; the global component covers the resources declared by !reads_of<self.pricing> *!.

State labels appear in two forms.
A \emph{single-state} predicate $\slabeled{S}{p}$ evaluates $p$ entirely in state $S$; the no-abort invariant uses this form because whether !pricing! aborts depends only on the pre-state of the call.
A \emph{two-state} predicate $\slabeled{S..}{q}$ takes $S$ as the pre-state of the !pricing! invocation and the resulting post-state as the second end; the output, monotonicity, and constant-product invariants use this form because the return value of !pricing! is only defined after execution from $S$.
More generally, $\slabeled{S_1..S_2}{q}$ names both ends explicitly; in the example, $\slabeled{S..}{q}$ fixes $S$ as the pre-state, with the post state unspecified.

The behavioral predicates !aborts_of!, !result_of!, and (elsewhere) !ensures_of! and !requires_of!, take a function value and its arguments and evaluate to the corresponding aspect of that call's behavior.
Without behavioral predicates, specifying a higher-order function would require inlining the callee's conditions at every call site, leading to a combinatorial explosion in specification size; behavioral predicates make reasoning about function values modular and abstract.
Using them in invariants as in the example creates (i)~a requirement for any code creating a pool to pass only compliant function values, and (ii)~a guarantee for any code calling a pool function that the requirements are met. Both conditions are verified by \MVP.

\Paragraph{Swapping}
The major functionality of the pool is to facilitate !swap! (exchange) of one asset with another.

% @formatter:off
\begin{MoveBox}
  fun swap(pool: &mut Pool, amt: u64): u64 {
    let out = (pool.pricing)
              (pool.reserve_x, pool.reserve_y, amt);
    pool.reserve_x += amt; pool.reserve_y -= out;
    out
  }
  spec swap {
    // Can only abort on reserve overflow.
    aborts_if pool.reserve_x + amt > MAX_U64;
    // Outcome is determined by pricing function.
    ensures result == old(result_of<pool.pricing>(
              pool.reserve_x, pool.reserve_y, amt));
    // Preserves product in pool.
    ensures pool.reserve_x * pool.reserve_y
         >= old(pool.reserve_x * pool.reserve_y);
  }
\end{MoveBox}
% @formatter:on

\noindent The specification of !swap! can be verified by \MVP because of the invariants guaranteed for the (dynamic) !pricing! function. In this case, we only illustrate the aborts condition and product preservation. For aborts, !pricing! itself is non-aborting, and (more subtly), |pool.reserve_y - out| cannot underflow because of the 'output never exceeds the output reserve' invariant, so only the overflow condition remains. (Notice \MVP would flag any aborts conditions not covered). For product preservation, the constraint can be derived as well.

\Paragraph{Pricing variants}
At the time a pool value is constructed, the invariants for the pricing function need to be verified.
Obviously this is more difficult than the application side, since now properties of functions need to be proven.
In reality, those proofs will often need custom solutions, but for the given example, \MVP can handle some cases gracefully.
First, we define a function which computes the so-called \emph{constant product}, which is a pricing function describing a direct swap without fees; however, rounding errors can lead to a loss in returned assets:

% @formatter:off
\begin{MoveBox}
  fun product(ri: u64, ro: u64, a: u64): u64 {
    let num = (ro as u128) * a;
    let den = (ri as u128) + a;
    if (den == 0) 0 else (num / den) as u64
  }
\end{MoveBox}
% @formatter:on

\noindent When constructing a pool using !Pool{.., pricing: product}!, \MVP is able to discharge all verification conditions.

The next example uses a pricing function which charges a configurable fee:

% @formatter:off
\begin{MoveBox}
  struct Fee has key {bps: u64}
  fun with_fee(addr: address,
               ri: u64, ro: u64, a: u64): u64 {
    let b = (a as u128)
       * (10000 - Fee[addr].bps as u128) / 10000;
    product(ri, ro, b)
  } spec with_fee {
    aborts_if !exists<Fee>(owner);
    aborts_if Fee[addr].bps > 10000;
    ..
  }
\end{MoveBox}
% @formatter:on

\noindent In order to construct a pricing function we need to curry the address for the fee resource and capture it in a closure, which can be done in Move as below:

% @formatter:off
\begin{MoveBox}
    let addr: address;
    Pool{
      ..,
      pricing: |ri, ro, a| with_fee(addr, ri, ro, a)
    }
\end{MoveBox}
% @formatter:on

\noindent \MVP detects that !with_fee! violates the requirement to not abort.
Guarding the definitions with checks and using a default fee fixes this and makes verification succeed.

% @formatter:off
\begin{MoveBox}
  fun with_fee_compliant(addr: address,
               ri: u64, ro: u64, a: u64): u64 {
    let bps =
      if (!exists<Fee>(addr) || Fee[addr].bps > 10000)
        500
      else
        Fee[addr].bps;
    let b =
        (a as u128) * (10000 - bps as u128) / 10000;
    product(ri, ro, b as u64)
  } spec with_fee_compliant {
    aborts_if false;
    ..
  }
\end{MoveBox}
% @formatter:on

Notice that currently, a user needs to provide complete aborts conditions of even simple functions as |with_fee|.
While \MVP in general allows to keep aborts conditions unspecified this does not work right now for aborts conditions in the case of function values.
We are looking at connecting specification inference in \MVP to the use case of compliance for function values; see also Sec.~\ref{sec:inference}.

The AMM example shows that non-trivial invariants can be specified for function values, and that the prover is capable of verifying them on the construction and caller sides, even though the invariants presented constitute non-linear arithmetic problems. While determining that the aborts contract is violated by |with_fee| appears easy (after all, one can directly see in the spec that the function aborts), there are more subtle things happening in the background -- for example, that |product| does not have any underflow on subtraction because of the |pricing| invariant that the return value does not exceed the reserve.

However, we have no illusions about the scalability of compliance-side verification.
In general, passing a function $f$ as an actual parameter $p$ requires proving $\forall x: \requiresof{p}{x} \implies \requiresof{f}{x}$, $\forall x: \abortsof{f}{x} \iff \abortsof{p}{x}$, and $\forall x,y: \ensuresof{f}{x, y} \implies \ensuresof{p}{x, y}$. For the examples provided here, this proof was possible. In general cases, we do not expect those proofs to be automated easily; rather, they will require proof hints and/or trusted assumptions. Nevertheless, on the application side of function values, such as inside the |swap| function, verification is not more complex than for first-order functions.

\SubSection{Filter-style Loops} \label{sec:find}

Behavioral predicates are not limited to struct invariants; they appear naturally in loop invariants and function specs for higher-order collection operations.
As a minimal illustration, consider a generic !find! which returns the first index of !v! whose element satisfies a predicate !pred!, or !len(v)! when no such element exists.
A spec-level helper !no_match_before! captures the recurring ``no earlier element has matched'' pattern:

% @formatter:off
\begin{MoveBox}
   fun find<T>(v: &vector<T>,
              pred: |&T|bool has copy + drop): u64 {
    let i = 0;
    let n = v.length();
    while (i < n) {
      if (pred(v[i])) return i;
      i += 1;
    } spec {
      invariant i <= n;
      invariant no_match_before(v, pred, i);
    };
    n
  }
  spec find {
    // opaque: caller side uses spec only
    pragma opaque;
    ensures result <= len(v);
    ensures no_match_before(v, pred, result);
    ensures result < len(v)
              ==> result_of<pred>(v[result])
  }
  spec fun no_match_before<T>(v: vector<T>,
                 pred: |&T|bool, end: u64): bool {
    forall j in 0..end: !result_of<pred>(v[j])
  }
\end{MoveBox}
% @formatter:on

\noindent !no_match_before! is an ordinary spec function, but its body calls the behavioral predicate !result_of<pred>! on each element --- a predicate-parametric abstraction is thus first-class in Move and can be reused across loop invariants and post-conditions.
The function's post-condition characterizes !result! entirely through the same evaluator: either !result! is a valid index whose call is the first to match, or !result == len(v)! and no element matches at all.
No concrete predicate is ever named in either the body or the spec; a caller instantiates !pred! with any function value, and the specification is discharged through its behavioral contract.

To illustrate the proof obligation at the caller, consider specializing !find! with an inline lambda whose behavioral contract is given by an inline !spec! clause.
As was discussed earlier for |with_fee|, functions provided as parameters need to have explicit specifications, therefore the lambda has a trivial spec block:

% @formatter:off
\begin{MoveBox}
  fun find_zero_lambda(v: &vector<u64>): u64 {
    find(v, |x| x == 0
             spec { ensures result == (x == 0) })
  }
  spec find_zero_lambda {
    ensures result <= len(v);
    ensures result < len(v) ==>
      v[result] == 0 &&
      (forall j in 0..result: v[j] != 0);
    ensures result == len(v) ==>
      (forall j in 0..len(v): v[j] != 0);
  }
\end{MoveBox}
% @formatter:on

\noindent Since !find! is opaque, \MVP reasons from !find!'s spec alone. (Without opaque, \MVP uses spec and body combined at caller side.)
At the match index, !find! guarantees $\resultof{pred}{v[result]}$, which --- by the lambda's spec --- reduces to $v[result] = 0$.
For earlier indices, !no_match_before! constrains $\resultof{pred}{v[j]}$ to be false for every index $j$ before $result$, which reduces to $v[j] \neq 0$.
\MVP discharges all three !ensures! clauses of !find_zero_lambda! without ever inspecting the lambda body.

Effectively, with the higher-order function |find|, we have verified a summary of this specific loop pattern which can now be reused many times on the caller side. Users calling functions like |find|, |map|, |reduce|, etc. will hardly need to write loop invariants themselves.


\SubSection{More About State Labels}
\label{sec:state-labels}

The Pool invariants in Sec.~\ref{sec:amm} used \emph{universally} quantified state labels: every state in which the pool exists must satisfy the pricing invariants.
State labels have a second, equally important use: \emph{existentially} quantifying over an intermediate state to express that a computation proceeds through a specific intermediate point.
Consider a higher-order function that sequences two stateful operations on a shared value:

% @formatter:off
\begin{MoveBox}
  fun followed_by(f: |&mut A|, g: |&mut A|, x: &mut A) {
    f(x); g(x)
  }
  spec followed_by {
    reads_of<f> *; reads_of<g> *;
    ensures exists S in *:
      (..S |~ ensures_of<f>(x)) &&
      (S.. |~ ensures_of<g>(x))
  }
\end{MoveBox}
% @formatter:on

\noindent The existentially quantified label !S! witnesses the intermediate state between the two calls.
Since no !modifies_of! is declared for !f! or !g!, they are assumed to have no global memory modification; the local component of the state at !S! is therefore the updated value of !x!, and the global component is unchanged.
The conjunct $\slabeled{..S}{\ensuresof{f}{x}}$ asserts that the postcondition of |f| holds from the pre-state of !followed_by! up to !S!; the conjunct $\slabeled{S..}{\ensuresof{g}{x}}$ asserts that the postcondition of !g! holds from !S! to the final post-state.
Crucially, the post-state of !f! automatically becomes the pre-state of !g!, without any explicit threading of intermediate values through the specification.

State labels are not restricted to higher-order functions.
When !f! and !g! are concrete named functions, the same pattern allows a wrapper to be specified by composing the behavioral predicates of its callees, avoiding duplication of their postconditions.


%%% Local Variables:
%%% mode: latex
%%% TeX-master: "main"
%%% End:
